\documentclass{article}
\usepackage{enumitem}
\usepackage{amsmath}

\begin{document}

\section*{Concept: Retrieval-Induced Forgetting (RIF)}

\subsection*{Question}

Maya is preparing for her French vocabulary test. She practices retrieving the French words for different categories of food---\emph{fruit}, \emph{vegetables}, and \emph{drinks}. She repeatedly quizzes herself only on the fruit words (e.g., ``apple'', ``orange'', ``pear''), ignoring the rest. Later that week, when she tries to recall all the vocabulary on the test, she remembers the practiced fruit words well, but she unexpectedly struggles to recall many of the vegetable words---even more so than the drink words that she never studied at all.

\vspace{1em}
\noindent
Which concept best explains Maya's pattern of forgetting?

\begin{enumerate}[label=\Alph*)]
    \item Proactive interference --- old learning interferes with new learning
    \item \textbf{Retrieval-induced forgetting} --- retrieving some items impairs recall of related, non-retrieved items
    \item Trace decay --- memory weakens simply with the passage of time
    \item Context-dependent memory --- recall depends on matching physical environments
\end{enumerate}

\subsection*{Correct Answer}

B) Retrieval-induced forgetting

\subsection*{Justification}

According to the retrieval-induced forgetting literature, practicing retrieval of certain items within a category impairs later recall of related, non-retrieved items from the same category. Maya repeatedly retrieves fruit-related vocabulary, which strengthens those associations but makes it harder to recall the related vegetable words. This pattern is characteristic of retrieval-induced forgetting.

\end{document}
